\scanpage{82}
\mspage{37}

\uline{Lemme} Soient $I$ un ens ordonné, $J,L$ deux parties de $I$. Conditions équivalentes

a) $1^{\mathrm o})$ $J,L$ sont des bandes, $2^{\mathrm o})$ tout élément de $J$ est majoré par un élément de $L$, tout élément de $L$ majore un élément de $J$, \deleted{$3^{\mathrm o})$ si $j\in J$, $l\in L$ sont tels que $l<j$, alors $j,l\in J\cap L$ !}
\margin{i.e. $J\prec L$}

b) $\exists$ une bande $I'$ de $I$, telle que $J$ soit un idéal de $I'$, $L$ un coidéal de $I'$, et que tout élément de $I'$ majore un élément de $J$ et soit majoré par un élément de $L$.

Sous ces conditions, \deleted{il existe \unsure{une plus}} \inserted{la bande $I'$ de b) est unique,} \unsure{et|c'est} la plus petite petite bande \inserted{$I'$} contenant $J$ et $L$ satisfait b), et est donnée par
\[
(67)\qquad I'=\{i\in I\mid\exists j\in J,l\in L\text{ avec }j\leq i\leq l\}
\]

\uuline{Dém} L'implication b) $\Rightarrow$ a) triviale, \inserted{ainsi que le fait que $I'$ est donné par (67),} prouvons a) $\Rightarrow$ b). Soit $I'$ donné par (67), on a $I'\supset J\cup L$ (par a) $2^{\mathrm o})$) $I'$ est une bande (trivial), et c'est évidemment la plus petite bande $\supset J\cup L$. Prouvons que $J$ est un idéal de $I'$ i.e. $j\in J$, $i\in I'$, $i\leq j\Rightarrow i\in J$.. Comme $i\in I'$, $\exists j_0\in J$ tel que $j_0\leq i$, donc $j_0\leq i\leq j$ avec $j_0,j\in J$, et comme $J$ est une bande \inserted{(par a $1^{\mathrm o})$)} on a $i\in J$. \unsure{On voit} directement que $L$ est un coidéal de $I'$. Soit $i\in I'$, \deleted{\unsure{de par}} définition de $I'$ $i$ majore un élément de $I$ et est majoré par un élément de $L$, qfd.
\mspage{81}
